% ============================================================
% DOCUMENTCLASS
% ============================================================
\documentclass[12pt,oneside]{book}

% ============================================================
% METADATOS
% ============================================================
\newcommand{\universidad}{Universidad de Buenos Aires (UBA)}
\newcommand{\materia}{Derecho Procesal Civil}
\newcommand{\titulo}{Sistemas de Solución de Conflictos}
\newcommand{\subtitulo}{Métodos Alternativos de Resolución de Conflictos}
\newcommand{\catedra}{Cátedra: Prof. Dr. Rojas}
\newcommand{\docenteclase}{Docente: Dra. Romina Moreno --- Adjunta de Cátedra (Comisiones 74-41 y 74-43)}
\newcommand{\autor}{Isaac Edgar Camacho Ocampo}
\newcommand{\anio}{2026}
\newcommand{\reflexionportada}{Comprender el conflicto es el primer paso para transformar la controversia en acuerdo.}

% ============================================================
% PAQUETES
% ============================================================
\usepackage[utf8]{inputenc}
\usepackage[T1]{fontenc}
\usepackage[spanish,es-nodecimaldot]{babel}

\usepackage[
    top=2.5cm,
    bottom=2.5cm,
    left=3cm,
    right=2.5cm,
    headheight=15pt
]{geometry}

\usepackage{amsmath}
\usepackage{amssymb}
\usepackage{mathtools}

\usepackage{graphicx}
\usepackage{float}
\usepackage{caption}
\usepackage{subcaption}

\usepackage{booktabs}
\usepackage{array}
\usepackage{longtable}
\usepackage{tabularx}

\usepackage{enumitem}

\usepackage{xcolor}
\usepackage[most]{tcolorbox}

\usepackage{fancyhdr}

\usepackage{ragged2e}

% ============================================================
% RUTAS DE IMÁGENES
% ============================================================
\graphicspath{
    {./img/}
    {./graficos/}
    {./figuras/}
}

% ============================================================
% HIPERVÍNCULOS Y METADATOS PDF
% ============================================================
\usepackage[
    colorlinks=true,
    linkcolor=blue!50!black,
    citecolor=green!40!black,
    urlcolor=blue!60!black,
    pdftitle={\titulo},
    pdfauthor={\autor},
    pdfsubject={Apuntes universitarios},
    bookmarks=true
]{hyperref}

% ============================================================
% CONFIGURACIÓN
% ============================================================
\setcounter{tocdepth}{2}

\setlength{\parindent}{0pt}
\setlength{\parskip}{0.5em}

\setlist[itemize]{
    topsep=0.3em,
    itemsep=0.2em
}

\setlist[enumerate]{
    topsep=0.3em,
    itemsep=0.2em
}

\clubpenalty=10000
\widowpenalty=10000

% ============================================================
% COMANDOS MATEMÁTICOS
% ============================================================
\newcommand{\abs}[1]{\left|#1\right|}
\newcommand{\norm}[1]{\left\lVert#1\right\rVert}
\newcommand{\vect}[1]{\mathbf{#1}}
\newcommand{\deriv}[2]{\frac{d#1}{d#2}}
\newcommand{\pderiv}[2]{\frac{\partial#1}{\partial#2}}

% ============================================================
% ENTORNOS / CAJAS ACADÉMICAS
% ============================================================
\newtcolorbox{definition}[1]{
    enhanced,
    breakable,
    title={Definición: #1},
    fonttitle=\bfseries,
    colback=blue!3,
    colframe=blue!50!black,
    boxrule=0.8pt,
    arc=2mm
}

\newtcolorbox{important}{
    enhanced,
    breakable,
    title={Importante},
    fonttitle=\bfseries,
    colback=yellow!8,
    colframe=orange!70!black,
    boxrule=0.8pt,
    arc=2mm
}

\newtcolorbox{example}[1]{
    enhanced,
    breakable,
    title={Ejemplo: #1},
    fonttitle=\bfseries,
    colback=green!4,
    colframe=green!40!black,
    boxrule=0.8pt,
    arc=2mm
}

\newtcolorbox{formula}[1]{
    enhanced,
    breakable,
    title={Fórmula / Regla: #1},
    fonttitle=\bfseries,
    colback=purple!4,
    colframe=purple!50!black,
    boxrule=0.8pt,
    arc=2mm
}

\newtcolorbox{exercise}[1]{
    enhanced,
    breakable,
    title={Ejercicio: #1},
    fonttitle=\bfseries,
    colback=gray!5,
    colframe=gray!60!black,
    boxrule=0.8pt,
    arc=2mm
}

\newtcolorbox{note}{
    enhanced,
    breakable,
    title={Nota},
    fonttitle=\bfseries,
    colback=white,
    colframe=gray!50,
    boxrule=0.6pt,
    arc=2mm
}

% ============================================================
% ENCABEZADOS Y PIES
% ============================================================
\pagestyle{fancy}
\fancyhf{}

\fancyhead[L]{\nouppercase{\leftmark}}
\fancyhead[R]{\materia}

\fancyfoot[C]{\thepage}

\renewcommand{\headrulewidth}{0.4pt}
\renewcommand{\footrulewidth}{0pt}

\fancypagestyle{plain}{
    \fancyhf{}
    \fancyfoot[C]{\thepage}
    \renewcommand{\headrulewidth}{0pt}
}

% ============================================================
% DOCUMENT
% ============================================================
\begin{document}

\frontmatter

% ------------------------------------------------------------
% PORTADA
% ------------------------------------------------------------
\begin{titlepage}

\thispagestyle{empty}

\begin{center}

\vspace*{1cm}

{\large \textsc{\universidad}}

\vspace{3cm}

{\Huge \textbf{\materia}}

\vspace{1cm}

{\LARGE \titulo}

\vspace{0.3cm}

{\large \subtitulo}

\vspace{0.6cm}

{\normalsize \catedra}

\vspace{0.2cm}

{\normalsize \docenteclase}

\vspace{0.8cm}

\textit{\reflexionportada}

\vspace{1.2cm}

\rule{0.70\textwidth}{0.5pt}

\vspace{0.5cm}

{\large Apuntes de estudio}

\vspace{0.5cm}

\rule{0.70\textwidth}{0.5pt}

\vfill

{\large \autor}

\vspace{0.3cm}

{\large \anio}

\end{center}

\vfill

\begin{flushleft}
\textit{Este es un modesto aporte a los estudiantes de ``\materia'' no pretende sustituir el material oficial de la materia.}
\end{flushleft}

\end{titlepage}

% ------------------------------------------------------------
% ÍNDICE
% ------------------------------------------------------------
\tableofcontents

\mainmatter

% ============================================================
% CAPÍTULO 1
% ============================================================
\chapter{Introducción a la Teoría del Conflicto}

\section{Presentación de la unidad}

La presente unidad aborda los sistemas de solución de conflictos, también denominados sistemas alternativos de solución de conflictos. El estudio parte de la teoría del conflicto, la cual reconoce dos fases de análisis: una fase estática y una fase dinámica. El desarrollo se concentra en el \textbf{análisis estático}, orientado a que el operador jurídico pueda identificar los elementos presentes cuando se encuentra frente a un conflicto. El \textbf{análisis dinámico}, vinculado a las escaladas del conflicto y a las herramientas para contenerlas, excede el alcance de esta unidad.

El análisis estático no integra la esencia del conflicto ---que es dinámico por naturaleza--- pero constituye una herramienta recomendable: permite mapear el conflicto, identificar de qué tipo se trata y, a partir de allí, intervenirlo y gestionarlo en su faceta dinámica.

\section{La formación del abogado como operador de conflictos}

Todo abogado tiene como materia prima el conflicto. Quien se desempeña como funcionario o empleado del Poder Judicial ---y en particular el juez--- tiene como tarea central resolver conflictos. Sin embargo, la materia \textit{Teoría del Conflicto} no reviste carácter obligatorio en los planes de estudio universitarios; en algunas asignaturas aparece únicamente como optativa.

\begin{important}
La ausencia de formación específica en teoría del conflicto conduce a que los operadores jurídicos aborden de un mismo modo conflictos de naturaleza muy distinta ---patrimoniales, de familia, sindicales, laborales, empresariales---, cuando cada tipo de conflicto requiere un mapeo propio, tanto desde su estática como desde su dinámica particular.
\end{important}

Frente al proceso por audiencias, la inmediación y las diversas posibilidades que ofrece el proceso judicial para gestionar conflictos, resulta necesario reflexionar sobre el tipo de formación que actualmente no están recibiendo los estudiantes de la carrera.

% ============================================================
% CAPÍTULO 2
% ============================================================
\chapter{Teoría del Conflicto: Análisis Estático}

\section{Definición de conflicto}

En torno a la definición de conflicto existen distintas perspectivas provenientes de la psicología, la economía y la filosofía. Dentro del ámbito jurídico se destacan las siguientes definiciones.

\begin{definition}{Cossio --- Escuela Egológica del Derecho}
El conflicto es una interferencia intersubjetiva de intereses. En la conducta humana, cuando existe interferencia entre dos personas en una relación intersubjetiva cuyos intereses son incompatibles, nace el conflicto.
\end{definition}

Carlos Cossio, referente de la Teoría Egológica del Derecho, sostenía que el derecho es conducta humana, definición desde la cual construye su concepto de conflicto.

Otra definición relevante corresponde a Remo Entelman, creador de la Teoría del Conflicto y autor de la obra \textit{Teoría del Conflicto}, en la que se desarrolla tanto el análisis estático como el dinámico del conflicto. En términos vinculados a la definición de Cossio, Entelman sostiene:

\begin{definition}{Entelman}
Hay conflicto cuando dos o más personas tienen objetivos incompatibles entre sí en una relación de interdependencia.
\end{definition}

De esta definición se desprenden tres elementos necesarios:
\begin{itemize}
    \item Dos o más personas.
    \item Objetivos incompatibles entre sí.
    \item Relación de interdependencia.
\end{itemize}

\begin{important}
El conflicto no requiere, para configurarse, la existencia de una norma que prohíba determinada conducta: dos personas pueden tener objetivos incompatibles aun cuando ambas desarrollen conductas permitidas por el derecho. Asimismo, el conflicto no debe identificarse necesariamente con la violencia, dado que puede no presentarse violencia en su desarrollo.
\end{important}

\section{Elementos del conflicto (análisis estático)}

Identificada la definición de conflicto, corresponde analizar los elementos que lo conforman desde una perspectiva estática. Se trata de un ejercicio que realiza el operador al momento de gestionar un conflicto, aun cuando en la realidad el conflicto sea dinámico.

Dos marcos teóricos resultan relevantes para este análisis:
\begin{itemize}
    \item \textbf{Teoría de la Comunicación}: estudia los elementos de la comunicación (emisor, receptor, mensaje, canal y código). La elección del canal a través del cual se transmite un mensaje ---por ejemplo, un mensaje de WhatsApp, un correo electrónico o una carta documento--- puede resultar determinante para evitar o propiciar escaladas del conflicto.
    \item \textbf{Teoría de la Información}: resalta la importancia de que el operador de conflictos recabe información y la organice de modo tal que resulte útil para la toma de decisiones. La información correctamente analizada genera la inteligencia necesaria para gestionar el conflicto.
\end{itemize}

\subsection{Los actores}

Los actores son los protagonistas del conflicto: quienes tienen una intención y una preocupación respecto del resultado, y pueden generar cambios en la dinámica del conflicto en pos de obtenerlo. En este contexto, el término \textit{actor} no se refiere al actor y al demandado del proceso judicial, sino a los protagonistas del conflicto concreto.

El actor del conflicto puede ser:
\begin{itemize}
    \item Individual.
    \item Colectivo, con liderazgo formal o informal, con organización o sin ella (por ejemplo, un grupo de personas que ocupa un inmueble no es asimilable a un sindicato, aunque ambos sean actores colectivos: cada supuesto debe analizarse específicamente).
\end{itemize}

Resulta relevante determinar si se encuentran presentes en la negociación todos los actores que deberían estarlo. Por ejemplo, en un conflicto de divorcio los actores directos son los cónyuges, pero la nueva pareja de una de las partes puede incidir considerablemente en el resultado ---en particular, en la división del régimen de comunidad de bienes---, por lo que puede ser necesario incorporarla, cuando menos, a la reunión entre el abogado y su cliente.

\begin{definition}{Conciencia del actor}
La conciencia es el producto de un acto intelectual mediante el cual el actor admite encontrarse, respecto de otro actor, en una situación en la que cree ---o ambos creen--- tener objetivos incompatibles.
\end{definition}

Trabajar sobre la conciencia del actor respecto de su situación de conflicto resulta fundamental. Existen personas que no tienen conciencia de encontrarse en un conflicto por no existir una norma que les prohíba determinada conducta. Por ejemplo, si ante un reclamo de mejoras laborales el abogado de la empresa señala que no existe obligación legal de conceder aumentos por no haberse firmado paritarias, puede inducir al empresario a creer que no está frente a un conflicto, cuando en rigor sí lo está.

\begin{important}
La falta de conciencia del conflicto puede llevar a que una de las partes no legitime a la otra en la mesa de negociación. Sin legitimar al adversario no es posible avanzar en la gestión de las distintas aristas que pueda presentar el conflicto ni alcanzar un acuerdo.
\end{important}

La percepción, por su parte, refiere al modo en que acceden a nuestro intelecto las opiniones, amenazas o riesgos existentes en la realidad, y que permiten anticipar cómo podría reaccionar el otro actor ante determinada conducta. Está estrechamente vinculada al conocimiento y a la información que se posee del otro actor.

\subsection{Los objetivos}

Los objetivos son aquello que persiguen los protagonistas del conflicto y que perciben como incompatible con los objetivos del otro actor. Se distinguen:
\begin{itemize}
    \item \textbf{Objetivos tangibles}: pueden verbalizarse y especificarse de modo concreto (por ejemplo: ``mi pretensión son 100 mil pesos'').
    \item \textbf{Objetivos ocultos}: no se explicitan, por razones estratégicas, de prestigio o de vergüenza, o porque el actor anticipa que, de verbalizarlo, la otra parte lo negará automáticamente en función del tipo de relación o de las características de su personalidad.
\end{itemize}

\subsection{Las relaciones de poder}

El poder se define, en la Teoría del Conflicto, como los recursos con que cuenta cada uno de los actores. Resulta relevante enumerar tanto los recursos de poder del propio cliente como los de la contraparte, a fin de determinar con qué herramientas se cuenta para gestionar la dinámica del conflicto.

Entre los recursos de poder se encuentran:
\begin{itemize}
    \item La posibilidad de generar una amenaza de sanción o una promesa de premio (por ejemplo, directivos de una empresa frente a sus trabajadores).
    \item El tiempo de espera disponible para alcanzar un acuerdo, que no reviste la misma relevancia para quien no necesita el dinero de inmediato que para quien lo requiere con urgencia.
    \item La paciencia.
    \item La alternativa al acuerdo.
    \item La posibilidad de efectuar la primera oferta.
    \item La información disponible sobre el caso y sus distintos elementos.
    \item El conocimiento que se posee del otro actor, especialmente relevante en conflictos de familia.
\end{itemize}

\subsection{Conflictos de suma cero y de suma variable}

Cuando se sostiene que todo lo que el otro actor gana implica necesariamente una pérdida propia, se adopta una mentalidad de \textbf{suma cero}, propia de los denominados \textit{conflictos puros}: siempre existiría un ganador y un perdedor. Por el contrario, en los \textit{conflictos impuros} rige la lógica de \textbf{suma variable}, en la cual existe una distribución de ganancias entre los actores.

Concebir un conflicto como de suma cero suele vincularse con la creencia de que no es posible generar soluciones que beneficien a ambas partes. Sin embargo, mediante un abordaje cooperativo y creativo pueden alcanzarse soluciones beneficiosas para todos los actores involucrados; en ello radica la importancia de la formación del operador de conflictos.

\subsection{Los terceros y las coaliciones}

Un último elemento del análisis estático consiste en evaluar qué terceros pueden intervenir en el conflicto o incidir en su resultado, así como la posibilidad de establecer coaliciones o alianzas estratégicas con alguno de ellos.

\begin{note}
Mapear el conflicto implica estudiar, para el caso concreto: quiénes son los actores (individuales o colectivos, su grado de cohesión y liderazgo, y si el actor puede fragmentarse o no); qué tipo de objetivos tienen (tangibles u ocultos); qué grado de conciencia poseen del conflicto; qué recursos de poder están disponibles; qué terceros pueden intervenir; y si es posible articular coaliciones o alianzas. Estos elementos, analizados desde la estática, inciden directamente en la toma de decisiones durante la fase dinámica del conflicto (manejo de escaladas y desescaladas mediante distintas técnicas de negociación).
\end{note}

% ============================================================
% CAPÍTULO 3
% ============================================================
\chapter{Sistemas de Resolución de Conflictos}

\section{Resolución y solución del conflicto}

El estudio de los sistemas de resolución de conflictos remite a un interrogante central: ¿cuál es la finalidad del derecho? Tradicionalmente se identifica con la paz social, es decir, con poner fin a los conflictos. Sin embargo, cabe preguntarse si la sentencia siempre pone fin al conflicto.

La función jurisdiccional culmina en un acto ---la sentencia--- que constituye la norma individual del caso y que determina un vencedor y un vencido. Ello no implica necesariamente que se haya puesto fin al conflicto. En un conflicto de familia, por ejemplo, lo que el juez determine puede no clausurar el conflicto, sino constituir un eslabón dentro de él, en la medida en que este tipo de conflictos presenta elementos adicionales ---emocionales, sentimentales, familiares--- que exceden la decisión judicial.

\begin{important}
Corresponde distinguir la \textbf{resolución} del conflicto ---formalmente contenida en la sentencia u otro acto procesal equivalente--- de la \textbf{solución} del conflicto, que puede vincularse con otros aspectos no alcanzados por dicho acto. De allí la importancia de que las partes, asistidas por operadores de conflictos, puedan darse su propia solución, la cual habitualmente resulta preferible a la que dispone el juez.
\end{important}

\section{Sistemas no adversariales}

Los sistemas alternativos de resolución de conflictos ---o de resolución alternativa de disputas--- se agrupan en dos grandes categorías, representables mediante una línea horizontal imaginaria: en un extremo se ubica la negociación directa entre las partes, sin intervención de terceros; en el otro extremo, el proceso judicial, con la máxima intervención de terceros para poner fin al conflicto.

Los \textbf{sistemas no adversariales} tienen como matriz a la \textbf{negociación}:
\begin{itemize}
    \item \textbf{Negociación directa} (entre las partes, autocomposición del conflicto): transacción directa entre las partes.
    \item \textbf{Negociación asistida} (interviene un tercero que colabora con las partes para posibilitar un acuerdo):
    \begin{itemize}
        \item Mediación.
        \item Conciliación.
        \item Otros sistemas (facilitadores, tercero neutral propio del sistema anglosajón, entre otros).
    \end{itemize}
\end{itemize}

\section{Sistemas adversariales}

En los \textbf{sistemas adversariales}, un tercero decide sobre el conflicto:
\begin{itemize}
    \item Proceso judicial.
    \item Arbitraje:
    \begin{itemize}
        \item Arbitraje de derecho.
        \item Arbitraje de amigables componedores.
        \item Pericia arbitral.
    \end{itemize}
\end{itemize}

\begin{note}
Suele hablarse de \textit{métodos alternativos} de resolución de conflictos. Corresponde, sin embargo, relativizar dicha denominación: calificar a estos métodos como ``alternativos'' presupone que existe otro método previo y ya previsto ---el proceso judicial--- respecto del cual los demás serían alternativos. Ello no se condice con la evolución histórica: la mediación, la conciliación y el arbitraje existen prácticamente desde los orígenes de la humanidad, como forma de resolución de conflictos en tribus, clanes y sectas, desde el derecho romano en adelante. El proceso judicial no fue, entonces, el método ``previo''.
\end{note}

\section{Comparación entre sistemas adversariales y no adversariales}

\begin{table}[H]
\centering
\begin{tabularx}{\textwidth}{@{}l X X@{}}
\toprule
\textbf{Criterio} & \textbf{No adversarial} & \textbf{Adversarial} \\
\midrule
Intervención de tercero & No siempre (solo en negociación asistida) & Siempre (juez o árbitro que decide) \\
Mecanismo procedimental & Rige el principio de informalidad & Está fijado por ley (proceso judicial) o por las partes (arbitraje) \\
Procedimiento probatorio & No es necesario (no hay que acreditar pretensiones) & Sí es necesario (el tercero evalúa la prueba) \\
Quién resuelve & Las propias partes se dan la solución & El juez o árbitro impone la decisión \\
Resultado & Acuerdo (con valor de sentencia si es homologado) & Sentencia o laudo arbitral \\
\bottomrule
\end{tabularx}
\caption{Diferencias entre sistemas no adversariales y adversariales}
\end{table}

% ============================================================
% CAPÍTULO 4
% ============================================================
\chapter{El Método de Harvard: Bases para la Negociación}

En el marco de los sistemas no adversariales, y tomando a la negociación como matriz, existen diversas escuelas y técnicas. Entre ellas se destaca el \textbf{Método de Harvard}, desarrollado por Roger Fisher y William Ury en la Escuela de Derecho de la Universidad de Harvard, quienes formularon sus bases en cuatro puntos.

\section{Separar a las personas del problema}

La primera base consiste en ser duro con el problema y suave con las personas. Con frecuencia, la negativa a escuchar al interlocutor o la agresión hacia la persona en razón del conflicto anulan la posibilidad de negociar. La persona no es el problema: el problema es uno, y las personas son distintas. En ocasiones, sustituir a los interlocutores ---por ejemplo, mediante la intervención de profesionales--- mejora la comunicación. Un abogado formado en teoría del conflicto y en negociación procurará mejorar la comunicación entre las partes para alcanzar un acuerdo, dando así cumplimiento a esta primera base.

\section{Concentrarse en los intereses, no en las posiciones}

Lo primero que manifiestan los actores de un conflicto son sus posiciones, que suelen presentarse como inamovibles e incompatibles. Sin embargo, las posiciones no reflejan necesariamente lo que las partes verdaderamente desean. El \textbf{interés} es aquello que subyace a la posición: el objetivo real perseguido, que puede resultar compatible con el de la otra parte.

\begin{example}{El conflicto de las naranjas}
Dos personas quieren la misma naranja; esa es su posición, y el conflicto parecería innegociable. Sin embargo, al indagar el interés real de cada parte, se advierte que una de ellas deseaba la naranja para comerla ---le interesaba el jugo o la pulpa--- mientras que la otra la deseaba para preparar dulce ---le interesaba la cáscara---. A partir de esa diferencia resulta posible alcanzar un acuerdo. Una cosa es la posición, y otra distinta es el interés.
\end{example}

\section{Generar opciones: el torbellino de ideas}

La tercera base consiste en generar diversas opciones antes de adoptar una decisión o elaborar una propuesta de acuerdo. Esta técnica se denomina \textit{torbellino de ideas}: permite que las partes ensayen una variedad de posibilidades de acuerdo, abandonando la hipótesis de que existe necesariamente un ganador y un perdedor, para explorar la posibilidad de alcanzar un acuerdo que beneficie a más de una parte. No todas las ideas generadas resultan factibles; lo relevante del ejercicio es descubrir posibilidades que inicialmente no se habían imaginado.

\section{Buscar criterios objetivos}

La cuarta base consiste en procurar criterios objetivos que sirvan de fundamento al acuerdo, en tanto contribuyen a convencer a las partes de que la solución propuesta resulta viable y razonable. Entre los ejemplos se encuentran la búsqueda de casos similares resueltos en la jurisprudencia, o la intervención de uno o varios tasadores para objetivar el valor de un inmueble o campo en disputa.

\begin{note}
Si las partes no logran darse su propia solución mediante estos criterios, el camino remanente es el sistema adversarial.
\end{note}

% ============================================================
% CAPÍTULO 5
% ============================================================
\chapter{La Mediación}

\section{Marco normativo}

La mediación se encuentra regulada en la Argentina por la Ley 26.589 y su decreto reglamentario N.\textsuperscript{o}~1.467/2011, normas incluidas entre las leyes complementarias del Código Procesal y que establecen la forma en que el mediador convoca la audiencia, su modo de designación y el procedimiento aplicable.

\begin{formula}{Ley 26.589 y Decreto Reglamentario 1.467/2011 --- Mediación previa y obligatoria}
En el ámbito nacional, la mediación está prevista como una instancia previa y obligatoria en los asuntos civiles o comerciales de contenido patrimonial, con las excepciones que la propia ley estipula.
\end{formula}

Agotada la instancia de mediación sin que se alcance un acuerdo, se labra un acta en la que se deja constancia de las partes intervinientes, del objeto de la mediación y de la falta de acuerdo. Dicha acta habilita la instancia judicial y constituye uno de los requisitos para promover la demanda.

\section{El mediador: requisitos y función}

Para desempeñarse como mediador se requiere:
\begin{itemize}
    \item Ser abogado matriculado.
    \item Contar con dos años de antigüedad en la matrícula.
    \item Reunir los requisitos adicionales exigidos para la habilitación correspondiente.
\end{itemize}
La habilitación es concentrada por el Ministerio de Justicia, en su carácter de dependencia del Poder Ejecutivo.

La función del mediador consiste en acercar a las partes y mejorar la comunicación entre ellas, de modo que sean las propias partes quienes alcancen un acuerdo respecto de sus diferencias. El mediador no debe proponer directamente fórmulas conciliatorias ni una posible solución al conflicto; su función se limita a facilitar que las partes encuentren su propia solución.

\section{Principio de confidencialidad}

En el sistema de mediación rige el principio de \textbf{confidencialidad}: lo manifestado, dicho o exhibido ---incluida la documentación o prueba presentada--- en el ámbito de la mediación no puede ser utilizado en un proceso judicial posterior.

\begin{important}
La confidencialidad permite que las partes negocien con la libertad necesaria para explorar las distintas aristas del conflicto. El mediador no puede ser convocado como testigo en un proceso judicial posterior para relatar lo manifestado por las partes durante la mediación.
\end{important}

\section{Mediación oficial y mediación privada}

En el ámbito civil y comercial, las partes pueden optar entre:
\begin{itemize}
    \item \textbf{Mediación oficial}: se solicita en la cámara del fuero correspondiente el sorteo de un mediador.
    \item \textbf{Mediación privada}: se elige un mediador de confianza adecuado al caso concreto, quien convoca a la otra parte mediante carta documento. Notificada la parte requerida, esta puede elegir entre cinco mediadores ofrecidos ---incluido el ya elegido por el requirente---. Si no elige ninguno, queda designado el mediador propuesto por el requirente; si elige uno de la lista, queda designado ese mediador.
\end{itemize}

\begin{note}
Esta posibilidad de elección no se encuentra prevista en la conciliación laboral previa.
\end{note}

% ============================================================
% CAPÍTULO 6
% ============================================================
\chapter{La Conciliación}

\section{Conciliación intra-procesal}

El artículo 360 del Código Procesal Civil y Comercial ---modificado por la Ley de Mediación 26.589--- prevé que los jueces, cuando adviertan que el conflicto podría resolverse mediante alguno de los métodos alternativos, pueden derivar el caso a mediación o solicitar la reapertura de esta, dentro de un plazo de treinta días. Asimismo, dicho artículo establece que, como una de las etapas de la audiencia preliminar ---la cual tiene lugar una vez transcurrida la etapa introductoria del proceso (demanda, contestación de demanda y prueba ofrecida)---, el propio juez invita a las partes a alcanzar una conciliación.

El artículo 36, inciso 2\textsuperscript{o}, del mismo Código faculta al juez para convocar a las partes, en cualquier momento del proceso y antes de dictarse sentencia ---sea en la etapa introductoria, probatoria o conclusional---, a los fines de intentar una conciliación.

A diferencia del mediador, el juez, en su función de conciliador, sí puede proponer fórmulas conciliatorias sobre la base de criterios objetivos, de la causa y de las cuestiones probatorias suscitadas, sin que ello importe un prejuzgamiento, en la medida en que el mecanismo está previsto precisamente para que así no ocurra.

\begin{table}[H]
\centering
\begin{tabularx}{\textwidth}{@{}l X X@{}}
\toprule
\textbf{Aspecto} & \textbf{Juez como conciliador} & \textbf{Mediador} \\
\midrule
Confidencialidad & No rige & Sí rige \\
Informalidad & No rige; deben cumplirse las formas del proceso & Sí rige \\
Dependencia institucional & Poder Judicial & Ministerio de Justicia (Poder Ejecutivo) \\
Proponer fórmulas conciliatorias & Sí puede & No puede \\
Nombramiento / requisitos & Es el juez de la causa & Abogado matriculado con dos años de antigüedad, habilitado por el Ministerio de Justicia \\
\bottomrule
\end{tabularx}
\caption{Diferencias entre el juez conciliador y el mediador}
\end{table}

\section{El acuerdo conciliatorio: artículo 309 CPCC}

El modo normal de terminación del proceso es la sentencia. No obstante, el Código Procesal prevé, entre los modos anormales de conclusión, la conciliación. Cuando en la audiencia preliminar ---u otro momento del proceso--- se alcanza un acuerdo, este adquiere la misma autoridad que una sentencia.

\begin{formula}{Art. 309 CPCC}
Los acuerdos conciliatorios celebrados por las partes ante el juez y homologados por este tendrán autoridad de cosa juzgada. Ante incumplimiento, se ejecutan por la vía de ejecución de sentencia.
\end{formula}

El mismo régimen se aplica al acuerdo celebrado en el marco de una mediación: el acta que recoge el acuerdo ---con sus deberes, obligaciones y previsiones ante incumplimiento--- tiene, para las partes, valor equivalente al de una sentencia. En caso de incumplimiento, dicho acuerdo se ejecuta por la vía de ejecución de sentencias (arts. 499 y siguientes; específicamente, para los acuerdos de mediación, art. 500 CPCC).

\section{SECLO: conciliación laboral obligatoria}

En los casos que tienen origen en el derecho del trabajo ---entre trabajadores y empleadores, diferencias salariales, cuestiones sindicales y demás materias propias del derecho laboral--- corresponde transitar la conciliación laboral obligatoria, prevista y regulada a través del SECLO (Servicio de Conciliación Laboral Obligatoria), que depende del Ministerio de Trabajo.

\begin{formula}{Ley 24.635}
Crea el SECLO y modifica la Ley de Procedimiento Laboral N.\textsuperscript{o}~18.345. El acta de cierre de la conciliación previa y obligatoria constituye un requisito para promover el juicio laboral.
\end{formula}

A diferencia del mediador, el conciliador laboral sí puede proponer fórmulas conciliatorias, de modo similar al juez. Asimismo, debe controlar el acuerdo celebrado ante sí, a los fines de su homologación por el Ministerio de Trabajo, en tanto se encuentran comprometidos el orden público y la protección de los derechos laborales; por ello, el acuerdo no puede celebrarse de cualquier modo.

\begin{note}
La materia civil y comercial ---y ciertas cuestiones de índole federal--- transita la mediación; la materia laboral transita la conciliación laboral obligatoria. Ambas constituyen instancias previas y obligatorias a la promoción del juicio.
\end{note}

% ============================================================
% CAPÍTULO 7
% ============================================================
\chapter{El Arbitraje}

El arbitraje se ubica entre los sistemas adversariales de resolución de conflictos: un tercero, el árbitro, interviene y decide sobre el conflicto. Se encuentra regulado en el Código Procesal Civil y Comercial a partir del artículo 736 y siguientes.

\section{Materia disponible y cláusula compromisoria}

Las partes pueden someter su conflicto a arbitraje únicamente cuando exista materia disponible, esto es, materia transable que no se encuentre afectada por el orden público.

\begin{important}
No puede someterse a arbitraje un conflicto de alimentos ni un conflicto de divorcio. Solo resulta arbitrable aquella materia disponible para las partes, sin afectación del orden público.
\end{important}

El origen del arbitraje es contractual: las partes acuerdan que, en caso de conflicto derivado del contrato, la diferencia será sometida a arbitraje. Dicha previsión se incorpora mediante la \textbf{cláusula compromisoria}: aquella cláusula por la cual las partes se comprometen a que su conflicto sea dirimido ante un arbitraje.

\section{Tipos de arbitraje}

\begin{itemize}
    \item \textbf{Arbitraje ad hoc}: se constituye al efecto para resolver un conflicto concreto. Las partes pueden designar a los árbitros ---por ejemplo, dos árbitros que a su vez eligen a un tercero, conformando un tribunal de tres miembros---.
    \item \textbf{Arbitraje institucional}: se somete el conflicto a un tribunal de arbitraje institucional ya existente, con reglamento propio y árbitros permanentes designados por concurso. Entre los ejemplos se encuentran el Tribunal de Arbitraje de la Bolsa de Comercio de Buenos Aires, el Tribunal de Arbitraje de la Bolsa de Cereales y el Tribunal de Arbitraje del Colegio de Abogados de San Isidro.
\end{itemize}

\section{Arbitraje de derecho}

\begin{formula}{Arbitraje de derecho}
El árbitro resuelve en los mismos términos en que podría hacerlo un juez, pero dentro de una jurisdicción privada, aplicando la norma y resolviendo conforme a derecho. Es, en general, más rápido y más económico en materia de tasas y gastos, y constituye el método de resolución de conflictos más utilizado a nivel internacional. La decisión se denomina \textbf{laudo} ---no sentencia--- y resulta equiparable a esta: ante incumplimiento, se ejecuta judicialmente por la vía de ejecución de sentencia (art. 499 y siguientes CPCC). Los árbitros carecen de facultades para ejecutar el laudo por sí mismos; para ello debe acudirse a la justicia.
\end{formula}

\section{Arbitraje de amigables componedores}

En esta modalidad, la resolución no se ajusta con la misma rigidez del derecho, sino que se resuelve conforme a la \textbf{equidad}. Los árbitros tienen en cuenta la norma y el derecho, pero aplican criterios de equidad. Según señala Barrios de Angelis, en este tipo de arbitraje se produce la ``dulcificación de la letra de la ley'': una forma de morigerar aquellas cuestiones fijadas taxativamente por la normativa, a fin de resolver conforme a la equidad del caso concreto.

\section{Pericia arbitral}

\begin{formula}{Pericia arbitral --- Art. 773 CPCC}
Prevista en el artículo 773 del Código Procesal Civil y Comercial, que remite al artículo 516. Procede cuando existe dificultad en las operaciones a realizar, en razón de que las liquidaciones resultan muy complejas. Para dirimir la disputa se designa a un experto en la materia, quien emite un dictamen a los fines de acercar posiciones o dirimir las diferencias.
\end{formula}

% TODO: las notas originales indican que el Dr. Rojas se plantea, en su trabajo sobre pericia arbitral, si esta puede constituir otro modo de conclusión del proceso, sin que el apunte precise una respuesta definitiva a ese interrogante.

% ============================================================
% CAPÍTULO 8
% ============================================================
\chapter{Referencias Bibliográficas y Normativas}

\section{Bibliografía recomendada}

\begin{itemize}
    \item Entelman, Remo. \textit{Teoría del Conflicto} (desarrolla el análisis estático y dinámico completo del conflicto).
    \item Calvo Soler, Raúl. \textit{Mapeo de Conflictos} (trabaja con proyecciones de casos concretos, combinando el análisis estático y dinámico).
    \item Rojas, Dr. Trabajos sobre arbitraje disponibles en \url{www.rojasomar.com.ar}, en particular: ``La pericia arbitral como otro modo de conclusión del proceso''.
    \item Fisher, R. y Ury, W. Método de Harvard sobre negociación.
\end{itemize}

\section{Normativa citada}

\begin{longtable}{@{}>{\raggedright\arraybackslash}p{0.28\textwidth} >{\raggedright\arraybackslash}X@{}}
\toprule
\textbf{Norma} & \textbf{Contenido / función} \\
\midrule
\endfirsthead
\toprule
\textbf{Norma} & \textbf{Contenido / función} \\
\midrule
\endhead
Ley 26.589 & Mediación prejudicial obligatoria; rige para materia civil y comercial en el ámbito nacional. \\
Decreto 1.467/2011 & Decreto reglamentario de la Ley 26.589. \\
Ley 24.635 & Crea el SECLO. Modifica la Ley de Procedimiento Laboral N.\textsuperscript{o}~18.345. Conciliación laboral obligatoria. \\
Art. 36 inc. 2\textsuperscript{o} CPCC & Facultad del juez de convocar a las partes para intentar una conciliación en cualquier momento del proceso. \\
Art. 360 CPCC & Audiencia preliminar; invitación a conciliación; derivación a mediación (modificado por Ley 26.589). \\
Art. 309 CPCC & Acuerdos conciliatorios con autoridad de cosa juzgada (modo anormal de terminación del proceso). \\
Arts. 499 y ss. CPCC & Ejecución de sentencia. Base para ejecutar laudos arbitrales y acuerdos de mediación. \\
Art. 500 CPCC & Específicamente, acuerdo de mediación: vía de ejecución. \\
Arts. 736 y ss. CPCC & Regulación del arbitraje en el Código Procesal Civil y Comercial. \\
Art. 773 CPCC & Pericia arbitral (remite al art. 516). \\
Art. 516 CPCC & Liquidaciones complejas (referenciado por el art. 773 sobre pericia arbitral). \\
\bottomrule
\end{longtable}

% ============================================================
% CAPÍTULO 9 --- CORNELL
% ============================================================
\chapter{Cuadro Cornell de Repaso}

El siguiente cuadro sintetiza, mediante el método Cornell, los conceptos centrales desarrollados en la unidad sobre sistemas de solución de conflictos.

\begin{longtable}{@{}>{\raggedright\arraybackslash}p{0.28\textwidth} >{\raggedright\arraybackslash}X@{}}
\toprule
\textbf{Preguntas / palabras clave} & \textbf{Notas / respuestas} \\
\midrule
\endfirsthead
\toprule
\textbf{Preguntas / palabras clave} & \textbf{Notas / respuestas} \\
\midrule
\endhead

\textbf{¿Qué es el conflicto? (Definición)} &
Cossio: interferencia intersubjetiva de intereses incompatibles. Entelman: existe conflicto cuando dos o más personas tienen objetivos incompatibles entre sí en una relación de interdependencia. Elementos: a) dos o más personas; b) objetivos incompatibles; c) relación de interdependencia. No requiere norma prohibitiva ni implica necesariamente violencia. \\

\textbf{¿Qué es el análisis estático del conflicto?} &
Es el mapeo previo del conflicto que realiza el operador antes de intervenir. Aunque el conflicto sea dinámico, este análisis permite identificar sus elementos para tomar decisiones adecuadas en la fase dinámica (manejo de escaladas y desescaladas). \\

\textbf{Elementos del análisis estático} &
1) Actores: individual o colectivo (con o sin liderazgo formal/informal, organizado o no), conciencia del conflicto y percepción. 2) Objetivos: tangibles (verbalizables) u ocultos (estratégicos, de prestigio, etc.). 3) Poder: recursos de cada actor (tiempo, paciencia, alternativas al acuerdo, información, conocimiento del otro, posibilidad de sanción o premio). 4) Terceros y posibilidad de coaliciones o alianzas. \\

\textbf{¿Qué es la conciencia del conflicto?} &
Acto intelectual por el cual un actor admite encontrarse en una situación de conflicto con otro. Sin conciencia, el actor no legitima al adversario, lo que bloquea toda gestión. Trabajar la conciencia resulta fundamental. \\

\textbf{Resolución vs. Solución} &
Resolución: acto formal (sentencia, laudo) que pone fin al proceso; siempre disponible en el sistema adversarial. Solución: satisfacción real de los intereses de las partes; no siempre la proporciona la sentencia (por ejemplo, en conflictos de familia con componente emocional). \\

\textbf{Suma cero vs. suma variable} &
Suma cero (conflictos puros): mentalidad de que todo lo que gana el otro lo pierdo yo; un ganador, un perdedor. Suma variable (conflictos impuros): distribución de ganancias; posibilidad de acuerdos cooperativos y creativos que beneficien a ambas partes. \\

\textbf{Sistemas no adversariales} &
Matriz: negociación. Directa (autocomposición): transacción entre partes. Asistida (heterocomposición): mediación, conciliación, facilitadores. Características: no siempre hay tercero; rige la informalidad; no hay procedimiento probatorio; las partes se dan su propia solución. \\

\textbf{Sistemas adversariales} &
Proceso judicial y arbitraje. Características: siempre hay un tercero que decide (juez o árbitro); mecanismo procedimental fijado (por ley o por las partes); hay procedimiento probatorio; el tercero evalúa la prueba y dicta resolución (sentencia o laudo). \\

\textbf{Método de Harvard (Fisher y Ury)} &
Cuatro bases: 1) separar personas del problema (duro con el problema, suave con la persona); 2) intereses, no posiciones (ir más allá de lo manifestado; ejemplo: el conflicto de las naranjas); 3) torbellino de ideas (generar múltiples opciones antes de decidir); 4) criterios objetivos (jurisprudencia similar, tasadores, peritos) para fundamentar el acuerdo. \\

\textbf{Mediación (Ley 26.589 / Dto. 1.467/2011)} &
Instancia previa y obligatoria en materia civil y comercial nacional. Designa: Ministerio de Justicia (Poder Ejecutivo). Función del mediador: acercar a las partes para que ellas encuentren su propia solución; no propone fórmulas. Rige la confidencialidad. Acuerdo con valor de sentencia, ejecutable por art. 500 CPCC. Modalidad oficial o privada (elección de mediador de confianza mediante carta documento). \\

\textbf{Conciliación judicial (arts. 360 y 36 inc. 2\textsuperscript{o} CPCC)} &
Art. 360: en la audiencia preliminar, el juez invita a las partes a una conciliación; también puede derivar a mediación (plazo: 30 días). Art. 36 inc. 2\textsuperscript{o}: en cualquier momento del proceso, antes de la sentencia, el juez puede convocar a una conciliación. El juez sí puede proponer fórmulas conciliatorias; no rige confidencialidad ni informalidad. Acuerdo logrado: art. 309 CPCC, autoridad de cosa juzgada; ejecución de sentencia si hay incumplimiento. \\

\textbf{Conciliación laboral (SECLO / Ley 24.635)} &
Instancia previa y obligatoria en materia laboral. Depende del Ministerio de Trabajo (Poder Ejecutivo). El conciliador sí puede proponer fórmulas conciliatorias. El acuerdo debe ser homologado por el Ministerio de Trabajo (control de orden público laboral). Norma: Ley 24.635 modifica la Ley 18.345. \\

\textbf{Arbitraje (arts. 736 y ss. CPCC)} &
Solo para materia disponible (sin afectar el orden público); no procede en alimentos, divorcio, etc. Nace de la cláusula compromisoria en el contrato. Tipos: ad hoc (creado para el caso) o institucional (Bolsa de Comercio, Bolsa de Cereales, Colegio de Abogados). Derecho: resuelve igual que un juez, aplica la norma, resultado = laudo (equiparado a sentencia, ejecutable por art. 499 CPCC). Amigables componedores: resuelve por equidad (``dulcificación de la letra de la ley'', Barrios de Angelis). Pericia arbitral (art. 773 CPCC, remite al art. 516): para liquidaciones complejas; especialistas emiten dictamen. \\

\textbf{¿Por qué ``alternativos'' entre comillas?} &
La mediación, la conciliación y el arbitraje existen desde el origen de la humanidad (tribus, clanes, Roma). El proceso judicial no fue ``previo''; hablar de métodos ``alternativos'' puede inducir a error histórico. \\

\midrule

\multicolumn{2}{p{\dimexpr\textwidth-2\tabcolsep}}{
\textbf{Resumen:} La unidad aborda los sistemas de resolución de conflictos desde dos grandes ejes: los no adversariales (negociación directa y asistida ---mediación, conciliación---) y los adversariales (proceso judicial y arbitraje). La Teoría del Conflicto (Cossio, Entelman) define el conflicto como la incompatibilidad de objetivos en una relación de interdependencia, sin requerir norma prohibitiva. El análisis estático (mapeo) identifica actores, objetivos, poder y terceros antes de gestionar el conflicto. La distinción entre resolución (acto formal: sentencia o laudo) y solución (satisfacción real de los intereses) es fundamental, en tanto la sentencia puede no poner fin a todos los conflictos, especialmente los de familia. El Método de Harvard (Fisher y Ury) propone cuatro bases: separar a las personas del problema, centrarse en los intereses y no en las posiciones, generar múltiples opciones mediante el torbellino de ideas, y emplear criterios objetivos. En cuanto al régimen normativo: la mediación (Ley 26.589) es previa y obligatoria en materia civil y comercial, rige la confidencialidad, y el acuerdo equivale a una sentencia (art. 500 CPCC). La conciliación judicial (arts. 36 inc. 2\textsuperscript{o} y 360 CPCC; acuerdo conforme al art. 309 CPCC) permite al juez proponer fórmulas. El SECLO (Ley 24.635) constituye la conciliación laboral previa y obligatoria. El arbitraje (arts. 736 y ss. CPCC) opera sobre materia disponible mediante cláusula compromisoria y puede ser de derecho, de amigables componedores o pericia arbitral (art. 773 CPCC).
} \\

\bottomrule
\end{longtable}

\end{document}
